\begin{tomtat}
	\begin{minipage}[t]{0.3\textwidth}
		\vspace{0pt}
		\begin{tomtat}
			\begin{center}
				\textbf{ĐỀ 25}
			\end{center}
			
		\end{tomtat}
	\end{minipage}
	\hfill
	\begin{minipage}[t]{0.5\textwidth}
		\vspace{0pt}
		\begin{tomtat}
			\begin{center}
				\includegraphics[scale=0.08]{images/25}
			\end{center}
		\end{tomtat}
	\end{minipage}
\end{tomtat}
\textbf{I. PHẦN TRẮC NGHIỆM}
\begin{ex}
	Kết quả rút gọn biểu thức $P=\sqrt{125}+\sqrt{5}$ là:
	\choice{$6\sqrt{5}$}{$-6\sqrt{5}$}{$5\sqrt{5}$}{$3\sqrt{5}$}
	\loigiai{
		Ta có:
		\begin{align*} P &= \sqrt{125} + \sqrt{5} \\ &= \sqrt{25 \cdot 5} + \sqrt{5} \\ &= 5\sqrt{5} + \sqrt{5} \\ &= (5+1)\sqrt{5} \\ &= 6\sqrt{5} \end{align*}
		\\ \textbf{Đáp Án A}
	}
\end{ex}

\begin{ex}
	Nghiệm của phương trình $4x - 1 = 15$ là:
	\choice{$\frac{7}{2}$}{$\frac{2}{7}$}{$4$}{$-4$}
	\loigiai{
		Ta có: $4x - 1 = 15$
		$\Leftrightarrow 4x = 15 + 1$
		$\Leftrightarrow 4x = 16$
		$\Leftrightarrow x = \frac{16}{4}$
		$\Leftrightarrow x = 4$.
		\\ \textbf{Đáp Án C}
	}
\end{ex}

\begin{ex}
	Cho tam giác ABC vuông tại A có $BC=8$cm, góc $B = 30^\circ$. Độ dài cạnh AC bằng:
	\choice{$16 \text{ cm}$}{$4 \text{ cm}$}{$4\sqrt{3} \text{ cm}$}{$5 \text{ cm}$}
	\loigiai{
		Trong tam giác ABC vuông tại A, ta có:
		$\sin B = \frac{AC}{BC}$
		$\sin 30^\circ = \frac{AC}{8}$
		$\frac{1}{2} = \frac{AC}{8}$
		$AC = 8 \cdot \frac{1}{2} = 4$ (cm).
		\\ \textbf{Đáp Án B}
	}
\end{ex}

\begin{ex}
	Cho hình nón có bán kính đáy $r=6$cm và chiều cao $h=8$cm. Diện tích xung quanh của hình nón là: (với $\pi \approx 3,14$)
	\choice{$150,72 \text{ cm}^2$}{$75,36 \text{ cm}^2$}{$94,2 \text{ cm}^2$}{$188,4 \text{ cm}^2$}
	\loigiai{
		Đầu tiên, ta cần tính độ dài đường sinh $l$ của hình nón.
		Áp dụng định lý Pitago trong tam giác vuông tạo bởi chiều cao, bán kính đáy và đường sinh:
		$l^2 = r^2 + h^2 = 6^2 + 8^2 = 36 + 64 = 100$.
		$l = \sqrt{100} = 10$ (cm).
		Diện tích xung quanh của hình nón được tính bằng công thức $S_{xq} = \pi r l$.
		$S_{xq} = \pi \cdot 6 \cdot 10 = 60\pi$ (cm$^2$).
		Với $\pi \approx 3,14$:
		$S_{xq} \approx 60 \cdot 3,14 = 188,4$ (cm$^2$).
		\\ \textbf{Đáp Án D}
	}
\end{ex}

\begin{ex}
	Một hình tròn có diện tích $S=144\pi \text{ cm}^2$. Tính bán kính của hình tròn đó là?
	\choice{$72 \text{ cm}$}{$12 \text{ cm}$}{$36 \text{ cm}$}{$14 \text{ cm}$}
	\loigiai{
		Diện tích hình tròn được tính bằng công thức $S = \pi R^2$.
		Theo đề bài, $S = 144\pi \text{ cm}^2$.
		$144\pi = \pi R^2$
		$R^2 = 144$
		$R = \sqrt{144} = 12$ (cm) (vì bán kính R > 0).
		\\ \textbf{Đáp Án B}
	}
\end{ex}

\begin{ex}
	Nghiệm của bất phương trình $3x+9>0$ là:
	\choice{$x>3$}{$x<-3$}{$x<3$}{$x>-3$}
	\loigiai{
		Ta có: $3x+9 > 0$
		$\Leftrightarrow 3x > -9$
		$\Leftrightarrow x > \frac{-9}{3}$
		$\Leftrightarrow x > -3$.
		\\ \textbf{Đáp Án D}
	}
\end{ex}
\begin{ex}
	Bảng sau ghi lại thời gian một bác sĩ khám cho một số bệnh nhân (đơn vị: phút)
	% Dữ liệu: 9,1; 7,7; 9,4; 9,1; 6,7; 5,9; 6,9; 6,0; 6,4; 8,7; 11,7; 5,8; 5,4; 6,4; 6,5; 12,3; 7,4; 10,0; 11,8; 5,3
	Chia số liệu thành 5 nhóm: [5; 6,5), [6,5; 8), [8; 9,5), [9,5; 11), [11; 12,5). Nhóm có tần số cao nhất là:
	\choice{$[5; 6,5)$}{$[6,5; 8)$}{$[9,5; 11)$}{$[11; 12,5)$}
	\loigiai{
		Dữ liệu thời gian khám bệnh (20 bệnh nhân):
		5,3; 5,4; 5,8; 5,9; 6,0; 6,4; 6,4; 6,5; 6,7; 6,9; 7,4; 7,7; 8,7; 9,1; 9,1; 9,4; 10,0; 11,7; 11,8; 12,3.
		
		Ta đếm tần số cho mỗi nhóm:
		\begin{itemize}
			\item Nhóm $[5; 6,5)$: Bao gồm các giá trị $x$ sao cho $5 \le x < 6,5$.
			Các giá trị: 5,3; 5,4; 5,8; 5,9; 6,0; 6,4; 6,4.
			Tần số: 7.
			\item Nhóm $[6,5; 8)$: Bao gồm các giá trị $x$ sao cho $6,5 \le x < 8$.
			Các giá trị: 6,5; 6,7; 6,9; 7,4; 7,7.
			Tần số: 5.
			\item Nhóm $[8; 9,5)$: Bao gồm các giá trị $x$ sao cho $8 \le x < 9,5$.
			Các giá trị: 8,7; 9,1; 9,1; 9,4.
			Tần số: 4.
			\item Nhóm $[9,5; 11)$: Bao gồm các giá trị $x$ sao cho $9,5 \le x < 11$.
			Các giá trị: 10,0.
			Tần số: 1.
			\item Nhóm $[11; 12,5)$: Bao gồm các giá trị $x$ sao cho $11 \le x < 12,5$.
			Các giá trị: 11,7; 11,8; 12,3.
			Tần số: 3.
		\end{itemize}
		Nhóm có tần số cao nhất là $[5; 6,5)$ với tần số là 7.
		\\ \textbf{Đáp Án A}
	}
\end{ex}

\begin{ex}
	Hộp sữa có dạng hình trụ với đường kính đáy là 12cm, chiều cao của hộp sữa là 18cm. Thể tích của hộp sữa bằng:
	\choice{$432\pi \text{ cm}^3$}{$648\pi \text{ cm}^3$}{$216\pi \text{ cm}^3$}{$2592\pi \text{ cm}^3$}
	\loigiai{
		Đường kính đáy của hình trụ là $d = 12$ cm.
		Bán kính đáy của hình trụ là $r = \frac{d}{2} = \frac{12}{2} = 6$ cm.
		Chiều cao của hình trụ là $h = 18$ cm.
		Thể tích của hình trụ được tính bằng công thức $V = \pi r^2 h$.
		$V = \pi \cdot (6)^2 \cdot 18$
		$V = \pi \cdot 36 \cdot 18$
		$V = 648\pi \text{ (cm}^3\text{)}$.
		\\ \textbf{Đáp Án B}
	}
\end{ex}

\newpage
\textbf{II. PHẦN TỰ LUẬN}
\begin{ex}
	Giải hệ phương trình:
	$$ \begin{cases} 2x-y=4 \\ 4x+6y=24 \end{cases} $$
	\loigiai{
		Ta có hệ phương trình:
		$$ \begin{cases} 2x-y=4 \quad &(1) \\ 4x+6y=24 \quad &(2) \end{cases} $$
		Từ phương trình (1), suy ra $y = 2x-4$.
		Thế $y = 2x-4$ vào phương trình (2):
		\begin{align*}
			4x + 6(2x-4) &= 24 \\
			4x + 12x - 24 &= 24 \\
			16x &= 24 + 24 \\
			16x &= 48 \\
			x &= \frac{48}{16} \\
			x &= 3.
		\end{align*}
		Thay $x=3$ vào $y = 2x-4$:
		\begin{align*}
			y &= 2(3) - 4 \\
			&= 6 - 4 \\
			&= 2.
		\end{align*}
		Vậy nghiệm của hệ phương trình là $(x; y) = (3; 2)$.
	}
\end{ex}

\begin{ex}
	Rút gọn các biểu thức sau:
	$$ B=\left(1+\frac{x+\sqrt{x}}{1+\sqrt{x}}\right)\left(1+\frac{x-\sqrt{x}}{1-\sqrt{x}}\right) $$
	với $0 \le x \neq 1$.
	\loigiai{
		Điều kiện: $0 \le x \neq 1$.
		Ta có:
		\begin{align*} B &= \left(1+\frac{\sqrt{x}(\sqrt{x}+1)}{1+\sqrt{x}}\right)\left(1+\frac{\sqrt{x}(\sqrt{x}-1)}{1-\sqrt{x}}\right) \\ &= \left(1+\sqrt{x}\right)\left(1-\frac{\sqrt{x}(1-\sqrt{x})}{1-\sqrt{x}}\right) \quad (\text{vì } 1-\sqrt{x} = -(\sqrt{x}-1)) \\ &= (1+\sqrt{x})(1-\sqrt{x}) \end{align*}
		Vậy $B = 1-x$.
	}
\end{ex}

\begin{ex}
	Gọi $x_1, x_2$ là hai nghiệm của phương trình bậc hai $x^2-14x+7=0$. Tính giá trị của biểu thức $x_1^2+x_2^2-2x_1x_2$.
	\loigiai{
		Phương trình $x^2-14x+7=0$ có $a=1, b=-14, c=7$.
		$\Delta' = b'^2 - ac = (-7)^2 - 1 \cdot 7 = 49 - 7 = 42 > 0$.
		Phương trình có hai nghiệm phân biệt $x_1, x_2$.
		Theo hệ thức Vi-ét, ta có:
		$$ \begin{cases} x_1 + x_2 = -\frac{b}{a} = -\frac{-14}{1} = 14 \\ x_1 x_2 = \frac{c}{a} = \frac{7}{1} = 7 \end{cases} $$
		Ta biến đổi biểu thức cần tính:
		\begin{align*} x_1^2+x_2^2-2x_1x_2 &= (x_1-x_2)^2 \\ &= (x_1+x_2)^2 - 4x_1x_2 \end{align*}
		Thay giá trị từ hệ thức Vi-ét vào:
		\begin{align*} (14)^2 - 4(7) &= 196 - 28 \\ &= 168 \end{align*}
		Vậy giá trị của biểu thức là 168.
	}
\end{ex}
\begin{ex}
	Hộp thứ nhất đựng 1 quả bóng trắng và 1 quả bóng đỏ. Hộp thứ hai đựng 1 quả bóng đỏ, 1 quả bóng vàng. Các quả bóng đều cùng kích thước và khối lượng. Lấy ngẫu nhiên mỗi hộp 1 quả bóng. Tính xác suất của biến cố sau: A: “2 quả bóng lấy ra có cùng màu”.
	\loigiai{
		Ký hiệu: Bóng trắng (T), bóng đỏ (Đ), bóng vàng (V).
		Hộp 1: $\{T_1; Đ_1\}$
		Hộp 2: $\{Đ_2; V_2\}$
		
		Lấy ngẫu nhiên mỗi hộp 1 quả bóng. Các kết quả có thể xảy ra (không gian mẫu $\Omega$):
		$\Omega = \{ (T_1; Đ_2); (T_1; V_2); (Đ_1; Đ_2); (Đ_1; V_2) \}$.\\
		Số phần tử của không gian mẫu là $n(\Omega) = 4$.\\
		
		Biến cố A: “2 quả bóng lấy ra có cùng màu”.
		Các kết quả thuận lợi cho biến cố A (hai quả bóng cùng màu đỏ):
		$A = \{ (Đ_1; Đ_2) \}$.\\
		Số phần tử của biến cố A là $n(A) = 1$.
		
		Xác suất của biến cố A là:
		$P(A) = \frac{n(A)}{n(\Omega)} = \frac{1}{4}$.
	}
\end{ex}
\begin{ex}
	Một khu vườn hình chữ nhật có chiều dài lớn hơn chiều rộng 3m và diện tích bằng 270m$^2$. Tìm chiều dài, chiều rộng của khu vườn.
	\loigiai{
		Gọi chiều rộng của khu vườn là $x$ (mét).
		Điều kiện: $x > 0$.
		Khi đó, chiều dài của khu vườn là $x+3$ (mét).
		
		Diện tích của khu vườn hình chữ nhật là 270m$^2$, nên ta có phương trình:
		\begin{align*} x(x+3) &= 270 \\ x^2 + 3x &= 270 \\ x^2 + 3x - 270 &= 0 \end{align*}
		Giải phương trình bậc hai $x^2 + 3x - 270 = 0$:
		Ta có $\Delta = b^2 - 4ac = 3^2 - 4(1)(-270) = 9 + 1080 = 1089$.
		$\sqrt{\Delta} = \sqrt{1089} = 33$.
		Phương trình có hai nghiệm phân biệt:
		$x_1 = \frac{-b + \sqrt{\Delta}}{2a} = \frac{-3 + 33}{2 \cdot 1} = \frac{30}{2} = 15$.
		$x_2 = \frac{-b - \sqrt{\Delta}}{2a} = \frac{-3 - 33}{2 \cdot 1} = \frac{-36}{2} = -18$.
		
		Đối chiếu với điều kiện $x > 0$:
		Nghiệm $x_1 = 15$ (thỏa mãn).
		Nghiệm $x_2 = -18$ (không thỏa mãn).
		
		Vậy, chiều rộng của khu vườn là 15 mét.
		Chiều dài của khu vườn là $x+3 = 15+3 = 18$ mét.
	}
\end{ex}
\begin{ex}
	Cho đường tròn tâm O, bán kính R. Từ một điểm M ở ngoài đường tròn, kẻ hai tiếp tuyến MA và MB với đường tròn (A, B là các tiếp điểm). Gọi E là trung điểm của đoạn thẳng MB, C là giao điểm của AE và (O) ($C \neq A$), H là giao điểm của AB và OM. Gọi D là giao điểm của MC và (O) ($D \neq C$).
	\begin{enumerate}
		\item[a)] Chứng minh: Các điểm A, M, B, O cùng thuộc một đường tròn.
		\item[b)] Chứng minh: $BE^2 = EC \cdot EA$ và tam giác ABD là tam giác cân.
	\end{enumerate}
	\loigiai{
	\begin{center}
		\includegraphics[scale=1.3]{hinh/25}
	\end{center}
	}
\end{ex}

\begin{ex}
	\begin{enumerate}
		\item[a)] Người ta dùng 100m rào để rào một mảnh vườn hình chữ nhật để thả gia súc. Biết một cạnh của hình chữ nhật là bức tường (không phải rào). Tính diện tích lớn nhất của mảnh vườn để có thể rào được.
		\item[b)] Cho các số thực $x, y, z > 2$. Tìm giá trị nhỏ nhất của biểu thức:
		$$ P = \frac{x}{\sqrt{y+z-4}} + \frac{y}{\sqrt{z+x-4}} + \frac{z}{\sqrt{x+y-4}} $$
	\end{enumerate}
\end{ex}

\setcounterref{ex}{0}
